First consider the simplex-vertex prepass. If \(\pi_{xy}=1\), nonnegativity and \(\sum_{x,y}\pi_{xy}=1\) imply that this unit cell is unique. Every iid increment equals \((x,y)\) almost surely. Hence the terminal strict-below counts are exactly \((Mx,My)\), and
\[
D_{M,k_0,k_1}(\pi)=\mathbf 1\{Mx<k_0\ \text{and}\ My<k_1\}.
\tag{1}
\]
Testing the four cells and the two inequalities in (1) takes constant exact time and constant workspace, including when a rank equals \(M+1\). Thus a vertex key is removed before any recurrence.

For a general key and \(0\le i\le M\), put
\[
R_i=\sum_{j=1}^i b_0(E_{j,0}),\qquad S_i=\sum_{j=1}^i b_1(E_{j,1}).
\]
We prove by induction that each retained entry satisfies \(d_i(r,s)=\Pr\{R_i=r,S_i=s\}\). This is the initialization at \(i=0\). Independence of draw \(i+1\) and the definition of the four cells give
\[
\Pr\{R_{i+1}=r,S_{i+1}=s\}
=\sum_{x,y\in\{0,1\}}\pi_{xy}\Pr\{R_i=r-x,S_i=s-y\}.
\tag{2}
\]
Negative indices have probability zero. Because both increments are nonnegative, a discarded state cannot return to the retained rectangle. Equation (2) is therefore exactly the asserted recurrence, and summing its terminal entries over \(r<k_0\), \(s<k_1\) proves the probability identity.

One scan of the table of \(R\) forms the four cells in \(O(|\mathscr E_0||\mathscr E_1|)\) operations. Each of the \(M\) updates visits at most \(k_0k_1\) states, and two rolling arrays use \(O(k_0k_1)\) memory. If, say, \(k_0=M+1\), then \(R_i<k_0\) is certain. Marginalizing the four cells over the first coordinate yields the analogous one-dimensional induction for \(S_i\), with \(O(Mk_1)\) operations and \(O(k_1)\) memory after tabulation. The other one-boundary case is symmetric, and with both ranks equal to \(M+1\) the joint event is certain. These arguments prove all proposed exactness and resource claims. The corner-oriented, multinomial-sum, and positive-corner refinements in the declared dependencies apply only after the constant-time vertex test and do not alter (2).